% Part: history
% Chapter: biographies 
% Section: kurt-goedel

\documentclass[../../../include/open-logic-section]{subfiles}

\begin{document}

\olfileid{his}{bio}{god}

\olsection{કુર્ત ગેડલ}

\olphoto{goedel-kurt}{કુર્ત ગેડલ}

કુર્ત ગેડલ (ઉચ્ચાર: ગેર-ડલ)નો જન્મ 28 એપ્રિલ
1906ના રોજ ઑસ્ટ્રો-હંગેરી સામ્રાજ્યના બ્રુનમાં
(આજે ચેક પ્રજાસત્તાકનું બ્રનો) થયો હતો. નાનકડા
કુર્ટલની જિજ્ઞાસુ અને તેજસ્વી પ્રકૃતિને કારણે
પરિવારમાં તેમને ઘણી વાર “Der kleine Herr Warum”
(નાનકડા શ્રી કેમ) કહેતા. પ્રાથમિક શાળાથી જ
તેઓ અભ્યાસમાં ઉત્તમ હતા; ગણિત એકમાત્ર વિષય
હતો જેમાં તેમને સર્વોચ્ચ ગુણ કરતાં ઓછા ગુણ
મળ્યા. નબળી તબિયતને કારણે ગેડલ ઘણી વાર
શાળામાં ગેરહાજર રહેતા અને તેમને શારીરિક
શિક્ષણમાંથી મુક્તિ આપવામાં આવી હતી. બાળપણમાં
તેમને સંધિવા તાવ હોવાનું નિદાન થયું હતું.
તબીબી તપાસમાં વિરુદ્ધ તારણ મળ્યું હોવા
છતાં, આ બીમારીથી પોતાના હૃદયને કાયમી
અસર થઈ છે એવું તેઓ જીવનભર માનતા રહ્યા.

ગેડલે 1924માં વિયેના યુનિવર્સિટીમાં અભ્યાસ
શરૂ કર્યો અને 1929માં ડૉક્ટરેટનો અભ્યાસ
પૂરો કર્યો. શરૂઆતમાં તેઓ ભૌતિકશાસ્ત્ર
ભણવા ઇચ્છતા હતા, પરંતુ ફિલસૂફ રુડોલ્ફ
કાર્નાપના પ્રભાવને કારણે પણ તેમની રુચિ
જલદી ગણિત અને ખાસ કરીને તર્કશાસ્ત્ર તરફ
વળી. હાન્સ હાનના માર્ગદર્શન હેઠળ લખેલા
તેમના નિબંધમાં તાદાત્મ્ય સાથેના પ્રથમ-ક્રમ
વિધેય તર્કશાસ્ત્રનું પૂર્ણતા પ્રમેય સાબિત
કરાયું હતું \citep{Godel1929}. માત્ર એક
વર્ષ પછી તેમણે પોતાનાં સૌથી પ્રખ્યાત
પરિણામો, પ્રથમ અને દ્વિતીય અપૂર્ણતા
પ્રમેયો, મેળવ્યાં (જે
\citealt{Godel1931}માં પ્રકાશિત થયાં).
વિયેનામાં રહેતાં ગેડલ વિજ્ઞાનલક્ષી
ફિલસૂફોના વિયેના વર્તુળમાં ખૂબ સક્રિય
હતા. તેમાં કાર્નાપ પણ હતા, જેમના
કાર્ય પર ગેડલનાં પરિણામોની ખાસ અસર પડી.

1938માં ગેડલે અડેલ નિમ્બુર્સ્કી સાથે
લગ્ન કર્યાં. તેમના માતાપિતા આથી ખુશ
ન હતા: અડેલ ગેડલ કરતાં છ વર્ષ મોટાં
અને અગાઉથી છૂટાછેડા લીધેલાં હતાં,
તેમજ રાત્રિક્લબમાં નૃત્યાંગના તરીકે
કામ કરતાં હતાં. છતાં સામાજિક દબાણની
ગેડલ પર અસર ન થઈ, અને તેમના
અવસાન સુધી બંનેનું લગ્નજીવન સુખી રહ્યું.

1938માં નાઝી જર્મનીએ ઑસ્ટ્રિયાને
જોડી લીધા પછી ગેડલ અને અડેલ
અમેરિકા સ્થળાંતરિત થયાં. ત્યાં ગેડલે
ન્યૂ જર્સીના પ્રિન્સ્ટનમાં આવેલી
ઇન્સ્ટિટ્યૂટ ફૉર ઍડવાન્સ્ડ સ્ટડીમાં
પદ સંભાળ્યું. અંતર્મુખી અને
અલગ સ્વભાવના હોવા છતાં પ્રિન્સ્ટનમાં
તેમનો સમય સહયોગપૂર્ણ અને ફળદાયી રહ્યો.
તેમણે ગણસિદ્ધાંત, તત્વજ્ઞાન અને
ભૌતિકશાસ્ત્ર પર નિબંધો પ્રકાશિત કર્યા.
ખાસ કરીને આ સંસ્થાના સહકર્મી
આલ્બર્ટ આઇન્સ્ટાઇન સાથે તેમની
ગાઢ મિત્રતા થઈ.

પછીનાં વર્ષોમાં ગેડલનું માનસિક સ્વાસ્થ્ય
કથળ્યું. 1977માં તેમની પત્નીને
હૉસ્પિટલમાં દાખલ કરવી પડી, તેથી
તેઓ ગેડલ માટે ભોજન બનાવી શકતાં
ન રહ્યાં.
ગેડલને જીવનભર માનસિક સ્વાસ્થ્યની
સમસ્યાઓ રહી હતી અને હવે તેઓ
વહેમથી ઘેરાઈ ગયા. ઝેર અપાશે
એવા ભારે ડરથી તેમણે ખાવાનું
બંધ કર્યું. 14 જાન્યુઆરી 1978ના
રોજ પ્રિન્સ્ટનમાં ભૂખમરાથી
તેમનું અવસાન થયું.

\begin{reading}
ગેડલનું સંપૂર્ણ જીવનચરિત્ર
\citet{Dawson1997}માં જુઓ. વધુ જીવનચરિત્રાત્મક
લેખો તથા તર્કશાસ્ત્ર અને તત્વજ્ઞાનમાં
ગેડલના પ્રદાન અંગેના નિબંધો માટે
\citet{Wang1990}, \citet{Baaz2011},
\citet{Takeuti2003} અને \citet{Sigmund2007} જુઓ.

ગેડલનો ડૉક્ટરેટ નિબંધ મૂળ જર્મન
ભાષામાં \citep{Godel1929} ઉપલબ્ધ છે.
અપૂર્ણતા પ્રમેયોનો મૂળ લેખ
\citep{Godel1931} છે. તેમનાં
પ્રકાશિત અને અપ્રકાશિત બધાં
લખાણો તથા પત્રવ્યવહારની પસંદગી
અંગ્રેજીમાં તેમના
\emph{Collected Papers}
\citet{Godel1986,Godel1990}માં ઉપલબ્ધ છે.

ગેડલના અપૂર્ણતા પ્રમેયોની
વિગતવાર ચર્ચા માટે
\citet{Smith2013} જુઓ.
તેમનાં પ્રમેયોની અનૌપચારિક
તત્વજ્ઞાનાત્મક ચર્ચા માટે
માર્ક લિન્સેનમાયરના પૉડકાસ્ટ
\citep{Linsenmayer2014}ને જુઓ.
\end{reading}

\end{document}
